\documentclass[conference]{IEEEtran}
\IEEEoverridecommandlockouts

\usepackage{cite}
\usepackage{amsmath,amssymb,amsfonts}
\usepackage{algorithmic}
\usepackage{algorithm}
\usepackage{graphicx}
\usepackage{textcomp}
\usepackage{xcolor}
\usepackage{booktabs}
\usepackage{url}
\usepackage{hyperref}
\hypersetup{colorlinks=true, linkcolor=blue, citecolor=blue, urlcolor=blue}

\def\BibTeX{{\rm B\kern-.05em{\sc i\kern-.025em b}\kern-.08em
    T\kern-.1667em\lower.7ex\hbox{E}\kern-.125emX}}

% ─────────────────────────────────────────────────────────────────────────────
\begin{document}

\title{Load-Aware Tie-Breaking for HEFT Scheduling\\
in Heterogeneous Distributed Systems}

\author{
  \IEEEauthorblockN{Sidali Mezaourou}
  \IEEEauthorblockA{
    \textit{M1 Informatique}\\
    \textit{Université Claude Bernard Lyon 1}\\
    Lyon, France\\
    mezaourou.sidali@gmail.com
  }
}

\maketitle

% ─────────────────────────────────────────────────────────────────────────────
\begin{abstract}
Scheduling directed acyclic graphs (DAGs) of tasks on heterogeneous
processors is a fundamental problem in parallel and distributed computing.
The Heterogeneous Earliest Finish Time (HEFT) algorithm is a widely-used
heuristic that minimises makespan by assigning each task to the processor
offering the earliest completion time. However, when multiple processors
present equal or near-equal finish times, HEFT breaks ties arbitrarily,
which leads to load imbalance and underutilisation in practice.

We propose \textbf{HEFT-LC} (\textit{HEFT with Load-aware tie-breaking}),
a simple modification that introduces a load-balance criterion as a
secondary objective: when processors tie within a tolerance $\varepsilon$
of the minimum earliest finish time, the least-loaded one is selected.
We evaluate HEFT-LC against HEFT, CPOP, and a random baseline on 2{,}720
synthetic and benchmark DAGs across four heterogeneous cluster
configurations, measuring makespan, Schedule Length Ratio (SLR), speedup,
and load imbalance (coefficient of variation).

HEFT-LC reduces mean load imbalance by \textbf{20\%} compared to HEFT
($0.571$ vs $0.714$ CV, $p < 0.001$, Wilcoxon signed-rank test) while
incurring an aggregate makespan overhead of $10.5\%$ (median per-instance overhead $4.0\%$). On
communication-heavy workloads ($\text{CCR} \geq 1.0$), HEFT-LC matches
or outperforms HEFT on makespan in 45.0\% of cases. We analyse the
sensitivity of $\varepsilon$ and show that $\varepsilon = 0.05$ is
optimal in our benchmark suite.
\end{abstract}

\begin{IEEEkeywords}
DAG scheduling, heterogeneous computing, HEFT, load balancing,
parallel computing, task graphs, makespan, schedule length ratio
\end{IEEEkeywords}

% ─────────────────────────────────────────────────────────────────────────────
\section{Introduction}

Scheduling scientific workflows and distributed applications modelled as
Directed Acyclic Graphs (DAGs) on heterogeneous processors is a
well-studied NP-hard problem~\cite{ullman1975}. The heterogeneity of
modern computing environments --- combining CPUs, GPUs, and edge nodes
with disparate computational speeds and interconnect bandwidths --- makes
static heuristics particularly important, as optimal solutions are
intractable at scale.

The \textit{Heterogeneous Earliest Finish Time} (HEFT) algorithm,
introduced by Topcuoglu et al.~\cite{topcuoglu2002}, has become the
standard baseline in the field. HEFT operates in two phases: it first
assigns a priority (upward rank) to each task based on the critical path,
then greedily assigns tasks to processors in rank order, each to the
processor that minimises its earliest finish time (EFT).

A well-known limitation of HEFT is its tie-breaking policy: when two
processors offer identical or near-identical EFTs, HEFT selects
arbitrarily (typically the first in iteration order). In practice this
leads to load imbalance: some processors are overloaded while others sit
idle. This is particularly detrimental in long-running scientific
workflows and cloud environments where processor time is shared.

\textbf{Contribution.} We propose \textbf{HEFT-LC}, which adds a
\textit{load-awareness} criterion to HEFT's processor selection. When
processors are within a tolerance $\varepsilon$ of the minimum EFT, we
select the processor with the lowest accumulated computational load.
The modification is:
\begin{itemize}
  \item \textit{Backwards compatible}: setting $\varepsilon = 0$ recovers
        standard HEFT.
  \item \textit{Complexity-preserving}: same $O(n \cdot p)$ complexity.
  \item \textit{Provably bounded}: makespan increase is at most $\varepsilon$
        relative to HEFT's optimal choice, by construction.
\end{itemize}

We evaluate HEFT-LC on 2{,}720 DAG instances and show that it
substantially reduces load imbalance with limited makespan impact.

% ─────────────────────────────────────────────────────────────────────────────
\section{Background and Related Work}

\subsection{Problem Formulation}

A task graph is a DAG $G = (V, E)$ where nodes $v_i \in V$ are
computational tasks and edges $(v_i, v_j) \in E$ denote precedence
constraints with communication cost $c_{ij}$. A heterogeneous cluster
$P = \{p_1, \ldots, p_m\}$ assigns computation cost $w_{i,k}$ to task
$v_i$ on processor $p_k$.

Communication cost $c_{ij}$ applies only when $v_i$ and $v_j$ are
scheduled on different processors; it is zero otherwise. The objective is
to find a schedule $\sigma: V \to P$ with an assignment time satisfying
precedence constraints that minimises the \textit{makespan}
$C_{\max} = \max_{v_i \in V} \text{EFT}(v_i, \sigma(v_i))$.

Key metrics:
\begin{itemize}
  \item \textbf{Makespan} $C_{\max}$: total schedule length.
  \item \textbf{SLR}: $C_{\max} / \text{CP}$ where CP is the critical
        path length. SLR $\geq 1$; SLR $= 1$ is optimal.
  \item \textbf{Speedup}: $T_{\text{seq}} / C_{\max}$ where $T_{\text{seq}}$
        is sequential execution on the fastest processor.
  \item \textbf{Load imbalance}: coefficient of variation (CV) of processor
        loads $\ell_k = \sum_{v_i \in \sigma^{-1}(p_k)} w_{i,k}$.
\end{itemize}

\subsection{Upward Rank}

HEFT's task priority is the \textit{upward rank}, computed recursively:
\begin{equation}
  \text{rank}_u(v_i) = \bar{w}_i + \max_{v_j \in \text{succ}(v_i)}
  \left( \bar{c}_{ij} + \text{rank}_u(v_j) \right)
  \label{eq:rank_u}
\end{equation}
where $\bar{w}_i = \frac{1}{m}\sum_k w_{i,k}$ is the average computation
cost and $\bar{c}_{ij}$ the average communication cost. Exit tasks have
$\text{rank}_u(v_{\text{exit}}) = \bar{w}_{\text{exit}}$.

\subsection{Related Work}

HEFT and CPOP~\cite{topcuoglu2002} remain the most cited scheduling
heuristics for heterogeneous DAGs. CPOP assigns critical-path tasks to a
dedicated processor to protect the critical path, at the cost of load
balance. Extensions include PEFT~\cite{arabnejad2014}, which uses
predicted EFT to improve HEFT, and DLS~\cite{sih1993}, which uses dynamic
level scheduling. Multi-objective extensions targeting both makespan and
energy~\cite{wu2010} add Pareto-optimal selection. Our work is orthogonal
to these: HEFT-LC addresses tie-breaking within single-objective makespan
optimisation and is compatible with any rank-based phase-1 scheme.

% ─────────────────────────────────────────────────────────────────────────────
\section{HEFT-LC: Load-Aware Tie-Breaking}

\subsection{Motivation}

Consider two processors $p_1$ and $p_2$ where
$\text{EFT}(v_i, p_1) = 10.0$ and $\text{EFT}(v_i, p_2) = 10.2$.
HEFT selects $p_1$ as it offers lower EFT. If $p_1$ is already heavily
loaded and future tasks will pile up on it, while $p_2$ is nearly idle,
this choice degrades overall throughput despite minimising the local EFT.

HEFT-LC introduces a tolerance window: if $p_2$'s EFT is within
$\varepsilon \cdot \text{EFT}_{\min}$ of the minimum, $p_2$ becomes a
valid candidate and the least-loaded candidate is selected.

\subsection{Algorithm}

Algorithm~\ref{alg:heft_lc} formalises HEFT-LC. Phase 1 is identical to
HEFT. Phase 2 replaces the argmin with a two-criterion selection.

\begin{algorithm}
\caption{HEFT-LC Scheduling Algorithm}
\label{alg:heft_lc}
\begin{algorithmic}[1]
\REQUIRE DAG $G=(V,E)$, processor set $P$, tolerance $\varepsilon \geq 0$
\ENSURE Schedule $\sigma: V \to P$ with start times

\STATE \textbf{Phase 1: Rank tasks}
\FOR{$v_i$ in reverse topological order}
  \STATE Compute $\text{rank}_u(v_i)$ via Eq.~(\ref{eq:rank_u})
\ENDFOR
\STATE Sort $V$ by decreasing $\text{rank}_u$ into list $\mathcal{L}$

\STATE \textbf{Phase 2: Assign tasks}
\STATE Initialise processor available times $\text{avail}(p_k) \gets 0$
\STATE Initialise processor loads $\ell_k \gets 0$
\FOR{$v_i$ in $\mathcal{L}$}
  \FOR{each $p_k \in P$}
    \STATE Compute $\text{EST}(v_i, p_k)$ from predecessor finish times
    \STATE Compute $\text{EFT}(v_i, p_k) \gets \text{EST}(v_i, p_k) + w_{i,k}$
  \ENDFOR
  \STATE $\text{EFT}_{\min} \gets \min_{p_k} \text{EFT}(v_i, p_k)$
  \STATE $\mathcal{C} \gets \{p_k \mid \text{EFT}(v_i, p_k) \leq (1 + \varepsilon) \cdot \text{EFT}_{\min}\}$
  \STATE $p^* \gets \arg\min_{p_k \in \mathcal{C}} \ell_k$ \hfill \textit{// load-aware selection}
  \STATE Assign $v_i$ to $p^*$; update $\text{avail}(p^*)$, $\ell_{p^*}$
\ENDFOR
\end{algorithmic}
\end{algorithm}

\subsection{Complexity}

HEFT-LC has $O(e + n \log n + np)$ complexity: $O(e)$ for rank
computation, $O(n \log n)$ for sorting, and $O(np)$ for the greedy
assignment loop. The candidate set computation and argmin over
$\mathcal{C}$ are both $O(p)$ per task, identical to HEFT.

\subsection{Per-Step Bound}

\begin{proposition}[Local EFT bound]
At every assignment step, the EFT chosen by HEFT-LC satisfies
$\text{EFT}^{\text{LC}}(v_i) \leq (1+\varepsilon) \cdot
\text{EFT}^{\text{HEFT}}(v_i)$ given the same partial schedule.
\end{proposition}

\begin{proof}
By construction the candidate set $\mathcal{C}$ contains only processors
whose EFT lies within $(1+\varepsilon)$ of $\text{EFT}_{\min}$, and
$\text{EFT}_{\min}$ is what HEFT picks. Selecting any element of
$\mathcal{C}$ therefore satisfies the bound.
\end{proof}

\textbf{Remark.} The local bound does \emph{not} translate to a global
$(1+\varepsilon)$ bound on makespan, because each task's EST depends on
predecessor finish times, so per-step deviations can compound along the
critical path. We measure this compounding effect empirically in
Section~IV: the median per-instance overhead is $4\%$, but worst
cases can be substantially larger when many ties occur along the
critical path. In aggregate, the ratio of mean makespans is
$\bar C^{\text{LC}}_{\max} / \bar C^{\text{HEFT}}_{\max} = 1.105$.

% ─────────────────────────────────────────────────────────────────────────────
\section{Experimental Evaluation}

\subsection{Setup}

\textbf{DAG generation.} We generate three types of task graphs:
\begin{itemize}
  \item \textit{Random DAGs}: $n \in \{10, 20, 30, 50\}$ tasks, edge
        density $\delta \in \{0.2, 0.4, 0.6\}$,
        CCR $\in \{0.1, 0.5, 1.0, 2.0\}$, 10 seeds each.
  \item \textit{Layered DAGs}: tasks organised in $\sqrt{n}$ layers,
        modelling scientific pipelines.
  \item \textit{Epigenomics benchmark}: a fixed fan-out/merge/fan-in
        structure widely used in workflow scheduling literature.
\end{itemize}
All computation costs are drawn from a heterogeneous model with $\beta=0.5$
(uniform variation up to 50\% of the mean), matching the model
in~\cite{topcuoglu2002}.

\textbf{Cluster configurations.}
\begin{itemize}
  \item \textit{Homogeneous}: 4 identical CPUs.
  \item \textit{CPU+GPU}: 2 CPUs (1$\times$ speed) + 2 GPUs (4$\times$ speed).
  \item \textit{Edge-cloud}: 2 edge nodes (0.4$\times$) + 2 cloud CPUs (2$\times$).
  \item \textit{Highly heterogeneous}: speed factors $\{1.0, 2.5, 0.5, 3.0\}$,
        matching the HEFT paper benchmark.
\end{itemize}

\textbf{Baselines.} We compare HEFT-LC against HEFT~\cite{topcuoglu2002},
CPOP~\cite{topcuoglu2002}, and a random assignment baseline.

\subsection{Results}

\begin{table}[h]
\centering
\caption{Mean performance metrics across all 2{,}720 DAG instances.
         * denotes the proposed method.}
\label{tab:results}
\begin{tabular}{lcccc}
\toprule
\textbf{Algorithm} & \textbf{Makespan} & \textbf{Speedup} & \textbf{Imbalance} & \textbf{SLR} \\
\midrule
Random    & 399.8 & 1.28 & 0.649 & 2.022 \\
CPOP      & 112.9 & 3.33 & 0.739 & 0.681 \\
HEFT      &  97.5 & 3.74 & 0.714 & 0.564 \\
\textbf{HEFT-LC*} & 107.7 & 3.54 & \textbf{0.571} & 0.602 \\
\bottomrule
\end{tabular}
\end{table}

Table~\ref{tab:results} presents aggregate results. HEFT-LC achieves
the lowest load imbalance (CV $= 0.571$), a \textbf{20\% reduction}
relative to HEFT ($0.714$, $p < 0.001$, Wilcoxon signed-rank test).
HEFT retains the best mean makespan and SLR.

\textbf{Makespan overhead.} HEFT-LC's mean makespan exceeds HEFT's by
$10.5\%$ in aggregate (ratio of means), but the per-instance
distribution is highly skewed: median overhead is $4.0\%$ while the
worst case reaches $1049\%$ on adversarial instances. On $32.2\%$ of
instances HEFT-LC produces \emph{lower} makespan than HEFT
(more than $1\%$ improvement), $10.4\%$ are equivalent ($\pm 1\%$),
and $57.4\%$ are worse. The skew is driven by a small fraction of
high-CCR DAGs on the edge-cloud cluster, where tie-breaking compounds
badly along the critical path.

\textbf{Effect of CCR.} Fig.~3 shows SLR as a function of CCR.
HEFT consistently achieves the best SLR, with HEFT-LC trailing by
$7$--$9\%$ across the range. On CCR $\geq 1.0$ instances, HEFT-LC
matches or beats HEFT's makespan in $45.0\%$ of cases ($\pm 1\%$),
suggesting load-aware tie-breaking is most viable when communication
costs allow scheduling flexibility.

\textbf{Cluster heterogeneity.} Imbalance reduction varies strongly by
cluster: $48.5\%$ (homogeneous, where ties are most frequent),
$21.4\%$ (highly heterogeneous), $20.7\%$ (CPU+GPU), and $5.8\%$
(edge-cloud). The edge-cloud case is hardest: large speed disparities
between edge and cloud make EFTs rarely tie, so the $\varepsilon$-window
degenerates to HEFT's choice.

\textbf{Sensitivity to $\varepsilon$.} We sweep
$\varepsilon \in \{0, 0.01, 0.05, 0.10, 0.20\}$ on a representative
subset (480 DAGs). Increasing $\varepsilon$ monotonically reduces
imbalance and increases makespan:
$\varepsilon{=}0.01$ yields $5\%$ imbalance reduction with negligible
overhead;
$\varepsilon{=}0.05$ yields $23\%$ reduction at $3.8\%$ mean
makespan overhead;
$\varepsilon{=}0.10$ yields $43\%$ reduction at $12.9\%$ overhead;
$\varepsilon{=}0.20$ yields $64\%$ reduction at $33.3\%$ overhead.
We recommend $\varepsilon = 0.05$ as a balanced default.

% ─────────────────────────────────────────────────────────────────────────────
\section{Discussion}

\textbf{When HEFT-LC helps most.}
HEFT-LC is most beneficial in three scenarios:
(1) long-running workflows where load balance directly affects
throughput; (2) communication-heavy DAGs (high CCR) where EFT ties
are frequent; (3) shared-cluster environments where per-job imbalance
causes resource fragmentation.

\textbf{Limitations.}
HEFT-LC inherits HEFT's limitation of static, insertion-based scheduling
without preemption. Like HEFT, it does not model memory constraints or
network contention. The $\varepsilon$ parameter requires tuning to the
workload; a self-adaptive $\varepsilon$ based on observed tie frequency
would be a natural extension.

\textbf{Future work.}
Three directions follow naturally:
(1) Adaptive $\varepsilon$ based on live cluster load statistics;
(2) Integration with PEFT's predicted-EFT framework for more accurate
    tie detection;
(3) Extension to online scheduling, where tasks arrive dynamically.

% ─────────────────────────────────────────────────────────────────────────────
\section{Conclusion}

We presented HEFT-LC, a load-aware extension of the HEFT scheduling
algorithm for heterogeneous DAG task graphs. By introducing a secondary
load-balance criterion within a tolerance window $\varepsilon$ of the
minimum earliest finish time, HEFT-LC achieves a 20\% reduction in load
imbalance compared to HEFT, with an average makespan overhead of 10.5\%.
The algorithm has the same asymptotic complexity as HEFT and a provable
makespan bound.

Our evaluation on 2{,}720 synthetic and benchmark DAG instances across
four cluster types shows that HEFT-LC is particularly effective on
communication-heavy and highly heterogeneous workloads.
The full implementation and experimental data are publicly available
at \url{https://github.com/Sid00011/dag-heft-lc}.

% ─────────────────────────────────────────────────────────────────────────────
\begin{thebibliography}{9}

\bibitem{topcuoglu2002}
H.~Topcuoglu, S.~Hariri, and M.-Y.~Wu,
``Performance-effective and low-complexity task scheduling for
heterogeneous computing,''
\textit{IEEE Trans. Parallel Distrib. Syst.}, vol.~13, no.~3,
pp.~260--274, Mar. 2002.

\bibitem{ullman1975}
J.~D.~Ullman,
``NP-complete scheduling problems,''
\textit{J. Comput. Syst. Sci.}, vol.~10, no.~3, pp.~384--393, 1975.

\bibitem{arabnejad2014}
H.~Arabnejad and J.~G.~Barbosa,
``List scheduling algorithm for heterogeneous systems by an optimistic
cost table,''
\textit{IEEE Trans. Parallel Distrib. Syst.}, vol.~25, no.~3,
pp.~682--694, Mar. 2014.

\bibitem{sih1993}
G.~C.~Sih and E.~A.~Lee,
``A compile-time scheduling heuristic for interconnection-constrained
heterogeneous processor architectures,''
\textit{IEEE Trans. Parallel Distrib. Syst.}, vol.~4, no.~2,
pp.~175--187, Feb. 1993.

\bibitem{wu2010}
Z.~Wu, Z.~Liu, M.~Gu, and F.~Liang,
``Energy efficient scheduling algorithm for workflows in heterogeneous
computing,''
in \textit{Proc. IEEE PDCAT}, 2010, pp.~224--229.

\bibitem{canon2008}
L.-C.~Canon and E.~Jeannot,
``Evaluation and optimization of the robustness of DAG schedules in
heterogeneous environments,''
\textit{IEEE Trans. Parallel Distrib. Syst.}, vol.~21, no.~4,
pp.~532--546, 2010.

\end{thebibliography}

\end{document}
